% ============================================================================ 2 · Contexto del proceso
\section{Contexto del proceso}
\label{sec:proceso}

\textbf{Qué entendimos.} La telemetría llega a la nube de forma continua, pero la degradación se identifica una vez
manifestada, en un service o cuando el cliente la reporta. El desafío se ubica en la posventa.

\textbf{Qué encontramos en los datos: el proceso reactivo está escrito en la telemetría.} La fecha registrada del evento
llega después del taller. En la primera entrega coincidía con la intervención (el cambio de aceite se concentraba en
±3 días). En la entrega v2, la fecha de los mismos autos cae una mediana de 14 días después, y en las dos semanas previas
ya aparecen el cambio de aceite y el ralentí largo de la regeneración forzada (1,5–1,8 veces la base). Después del evento,
además, el uso cambia: la fracción de viajes bajo régimen térmico baja en el 72\,\% de los autos.

\textbf{Qué hicimos.} Si el modelo viera el taller, haría la detección reactiva que Ford ya tiene. Por eso la referencia
del evento es la \textbf{fecha registrada menos 21 días}, el modelo nunca ve los \textbf{500 km previos} a esa referencia
(§\ref{sec:panel}) y nada posterior al evento entra a ninguna ventana. En la posventa, la alerta entra a un circuito que
termina en una acción (figura~\ref{fig:circuito}): con un hábito para nombrar, un consejo al conductor en la app; sin
hábito, un llamado del concesionario; si el riesgo sigue alto dos revisiones después, un escalamiento. El gerente de
posventa recibe cada lunes una bandeja priorizada y redactada, lista para aprobar.

\begin{figure}[!htbp]
\centering
\begin{tikzpicture}[
  node distance=4mm and 3.5mm,
  box/.style={draw=fordblue!70,fill=softbg,rounded corners=2pt,align=center,font=\sffamily\scriptsize,
    minimum height=10mm,text width=21mm,inner sep=2pt},
  pol/.style={box,draw=accent2!80,fill=warnbg},
  arr/.style={-{Stealth[length=2mm]},draw=inksec,thick}]
\node[box] (t) {\textbf{1 · Telemetría}\\viajes y señales del filtro};
\node[box,right=of t] (s) {\textbf{2 · Puntaje GRU}\\cada 500 km};
\node[box,right=of s] (a) {\textbf{3 · Alerta}\\2 cortes sobre el umbral de Ford};
\node[box,right=of a] (w) {\textbf{4 · Por qué}\\hábitos contra los sanos del mercado};
\node[pol,right=of w] (p) {\textbf{5 · Política}\\determinista, la define Ford};
\node[box,below=5mm of p] (g) {\textbf{6 · Agentes}\\redactan; un verificador controla};
\node[box,left=of g,text width=30mm] (u) {\textbf{Destinatarios}\\conductor · concesionario · posventa};
\node[box,left=of u,text width=30mm] (r) {\textbf{7 · Retorno del taller}\\monitoreo y reentrenamiento};
\draw[arr] (t) -- (s); \draw[arr] (s) -- (a); \draw[arr] (a) -- (w); \draw[arr] (w) -- (p);
\draw[arr] (p) -- (g); \draw[arr] (g) -- (u); \draw[arr] (u) -- (r); \draw[arr] (r) -| (t);
\end{tikzpicture}
\caption{El circuito de posventa: el modelo decide, el agente comunica y el código cuida lo que se dice.}
\label{fig:circuito}
\fuente{\repo{docs/\allowbreak{}pitch/\allowbreak{}presentacion-contenido.md} bloque 4; \repo{feat/\allowbreak{}demo-gru:PRODUCT.md}.}
\end{figure}
